\begin{helper}
Let \(C\) be a base event and \(T\) a target event in the finite two-bite
sample space, with \(0\le p\le1\), and let probabilities be computed by
\(\noderef{TwoBiteEventProbability}\).  Suppose \(T\cap C\) is covered by a
finite family of cells
\[
G_{b,s,t},
\]
where \(b\) indexes complete first-color traces, \(s\) indexes later support
choices, and \(t\) indexes prefix transcripts.  Suppose each cell has a
nonnegative assigned weight \(w_{b,s,t}\), and that
\[
\Pr(T\cap C\cap G_{b,s,t})\le w_{b,s,t}
\]
for every triple \((b,s,t)\).  If \(M\ge0\) and
\[
\sum_{b,s,t} w_{b,s,t}\le M,
\]
then
\[
\Pr(T\cap C)\le M.
\]
\end{helper}
\begin{proof}
By \(\noderef{TwoBiteSampleWeightNonnegative}\), the assumptions
\(0\le p\le1\) make every sample weight nonnegative.  Thus every probability
below is a finite sum of nonnegative weights.  For every sample
\(\omega\in T\cap C\), the covering hypothesis supplies at least one triple
\((b,s,t)\) with \(\omega\in G_{b,s,t}\).  Therefore
\[
T\cap C\subseteq \bigcup_{b,s,t}(T\cap C\cap G_{b,s,t}).
\]
Finite subadditivity for nonnegative weights gives
\[
\Pr(T\cap C)
\le
\sum_{b,s,t}\Pr(T\cap C\cap G_{b,s,t}).
\]
Apply the assumed cell estimate to every summand:
\[
\sum_{b,s,t}\Pr(T\cap C\cap G_{b,s,t})
\le
\sum_{b,s,t} w_{b,s,t}.
\]
The final assumed weight-sum bound gives
\[
\Pr(T\cap C)\le M.
\]
No factor counting the number of complete traces appears: the traces and
prefix transcripts enter only through their actual cylinder weights, and the
hypothesis requires those weights themselves to sum to \(M\).
\end{proof}
